\begingroup
\small
\setlength{\tabcolsep}{4pt}
\renewcommand{\arraystretch}{1.15}
\begin{longtable}{>{\raggedright\arraybackslash}p{3.0cm}>{\raggedright\arraybackslash}p{5.6cm}>{\raggedright\arraybackslash}p{5.9cm}}
\caption{Synthesis of prior contributions and implications for TDRP and Masar.}\label{tab:lr_synthesis}\\
\toprule
\textbf{Decision area} & \textbf{Established contribution} & \textbf{Implication for the proposed research} \\
\midrule
\endfirsthead
\multicolumn{3}{l}{Table \thetable\ continued}\\
\toprule
\textbf{Decision area} & \textbf{Established contribution} & \textbf{Implication for the proposed research} \\
\midrule
\endhead
\bottomrule
\endfoot
Continuity requirements & Contextual BIA, recovery objectives, and resources in government guidance \cite{nist2010contingency,dga2025guidelineen}. & Preserve the organizational meaning of parameters; do not claim that software alone demonstrates compliance. \\
Temporal assessment & Time-indexed impact matrices in AMM BIA and configurable temporal calculations in product documentation \cite{amm2023bia,globalsuite2025prebia,servicenow2025rto}. & Specify the exact rule semantics, boundaries, missing-data behavior, and configuration used by TDRP. \\
Criticality aggregation & Weighted, fuzzy, and probabilistic assessment of business-derived asset importance \cite{assetcriticality2023}. & Explain why a chosen aggregation policy fits the decision; do not assume a simple rule is universally superior. \\
Recovery execution & Resource-constrained operating trajectories and dependency-aware restoration scheduling \cite{sahebjamnia2015integrated,elshekeil2017precesion}. & Keep service ranking distinct from operational sequencing and capacity allocation. \\
Supporting dependencies & Business-derived asset criticality and process-level impact propagation \cite{quinn2025bia,nist2018criticality,raptaki2025propagation}. & Define shared-dependency inheritance explicitly; distinguish classification from causal propagation. \\
Artifact evaluation & Design-science methods and evaluation strategies separating properties and settings \cite{hevner2004design,peffers2007dsrm,venable2016feds}. & Evaluate method, implementation, and practitioner interpretation separately; bound operational claims. \\
\end{longtable}
\endgroup
